\chapter*{Preface}
\addcontentsline{toc}{chapter}{Preface}

This book, is designed to guide students through the foundational concepts of computer programming using a hands-on, project-based approach. By consistently working with a single dataset—the Iris dataset—we aim to build a cohesive understanding of data manipulation, algorithmic thinking, and the basics of machine learning.

\section*{Goal and Strategies}

The primary goal of this tutorial series is to transition students from basic syntax to building functional, rule-based classifiers. The strategy employed is one of incremental complexity:
\begin{itemize}
    \item \textbf{Continuity:} We use the same dataset (Iris) across all sessions to reduce cognitive load and allow focus on new programming concepts.
    \item \textbf{Scaffolding:} Each session builds directly on the previous one, starting with simple variables and progressing to functions, loops, and modular design.
    \item \textbf{Practical Application:} Theoretical concepts are immediately applied to solve concrete problems, reinforcing learning through doing.
    \item To avoid installation issues, we will use Github Codespaces for all coding exercises, ensuring that students can focus on learning rather than troubleshooting their development environment.
\end{itemize}

\section*{Intended Audience}

This material is written for:
\begin{itemize}
    \item Undergraduate students in computer science, data science, or related engineering fields.
    \item Beginners with little to no prior programming experience who want a structured introduction to Python.
    \item Learners interested in the intersection of programming and introductory data analysis.
\end{itemize}

\section*{Acknowledgements}

This tutorial series is the first edition,and it is a work in progress. Therefore, I welcome feedback and suggestions for improvement. If you have any comments or ideas on how to enhance the content, please feel free to reach out to me through the project's GitHub repository.


\section*{Repository}
All code examples and exercises for this tutorial series are available in the GitHub repository: \url{https://github.com/balandongiv/computer_programming_codespaces}. You should fork the repository and use it as a reference throughout the sessions. Each session's code is organized in separate folders for easy navigation.
